\section{USAJMO 2013}
        \begin{enumerate}
        	\item (\textit{USAJMO 2013}) Are there integers $a$ and $b$ such that $a^5b+3$ and $ab^5+3$ are both perfect cubes of integers?


 

	\item (\textit{USAJMO 2013}) Each cell of an $m\times n$ board is filled with some nonnegative integer.  Two numbers in the filling are said to be \textit{adjacent} if their cells share a common side.  (Note that two numbers in cells that share only a corner are not adjacent).  The filling is called a \textit{garden} if it satisfies the following two conditions:


 (i) The difference between any two adjacent numbers is either $0$ or $1$.


 (ii) If a number is less than or equal to all of its adjacent numbers, then it is equal to $0$.


 Determine the number of distinct gardens in terms of $m$ and $n$.


 

	\item (\textit{USAJMO 2013}) In triangle $ABC$, points $P,Q,R$ lie on sides $BC,CA,AB$ respectively.  Let $\omega_A$, $\omega_B$, $\omega_C$ denote the circumcircles of triangles $AQR$, $BRP$, $CPQ$, respectively.  Given the fact that segment $AP$ intersects $\omega_A$, $\omega_B$, $\omega_C$ again at $X,Y,Z$ respectively, prove that $YX/XZ=BP/PC$.


 

	\item (\textit{USAJMO 2013}) Let $f(n)$ be the number of ways to write $n$ as a sum of powers of $2$, where we keep track of the order of the summation.  For example, $f(4)=6$ because $4$ can be written as $4$, $2+2$, $2+1+1$, $1+2+1$, $1+1+2$, and $1+1+1+1$.  Find the smallest $n$ greater than $2013$ for which $f(n)$ is odd.


 

	\item (\textit{USAJMO 2013}) Quadrilateral $XABY$ is inscribed in the semicircle $\omega$ with diameter $XY$.  Segments $AY$ and $BX$ meet at $P$.  Point $Z$ is the foot of the perpendicular from $P$ to line $XY$.  Point $C$ lies on $\omega$ such that line $XC$ is perpendicular to line $AZ$.  Let $Q$ be the intersection of segments $AY$ and $XC$.  Prove that 
 \[\dfrac{BY}{XP}+\dfrac{CY}{XQ}=\dfrac{AY}{AX}.\]


 

	\item (\textit{USAJMO 2013}) Find all real numbers $x,y,z\geq 1$ satisfying 
 \[\min(\sqrt{x+xyz},\sqrt{y+xyz},\sqrt{z+xyz})=\sqrt{x-1}+\sqrt{y-1}+\sqrt{z-1}.\]


 

\end{enumerate}